\chapter{Introducción}

\section{Motivación}

El sistema QCAL (Quantum Coherence Algebraic Logic) emerge de la observación
fundamental de que la coherencia cuántica en espacios adélicos produce una
frecuencia natural de resonancia: $f_0 = 141.7001$ Hz.

Esta frecuencia no es arbitraria. Surge como punto fijo del flujo de
renormalización que gobierna la dinámica del operador de Dirac adélico
$\mathfrak{D}$ sobre el espacio de Hilbert $L^2(\mathbb{A})$, donde
$\mathbb{A}$ denota el anillo de adeles.

\section{Arquitectura del Sistema}

El sistema QCAL se estructura en cinco pilares formales:

\begin{enumerate}
  \item \textbf{Cuantización Espectral:} El espectro de $\mathfrak{D}$ es real y discreto,
        demostrado desde la topología adélica.
  \item \textbf{Emergencia de $f_\text{eff}$:} $f_0 = 141.7001$ Hz es el mínimo
        global de la Acción Noética de Ginzburg-Landau.
  \item \textbf{Equivalencia QCAL-RH:} La hipótesis de Riemann es equivalente a la
        conjetura espectral del sistema QCAL.
  \item \textbf{Seguridad Lógica:} El sistema impide replay y doble gasto por
        construcción formal.
  \item \textbf{Unificación del Acoplo:} Cuatro caminos independientes convergen a
        $\Delta_c = 0.1054947116$ Hz.
\end{enumerate}
